\section{Determinantes da Conversão da Vegetação Nativa em Mato Grosso (2001-2024)}
\label{literatura}

Compreender a relação entre expansão agropecuária e conversão da vegetação nativa envolve discutir em quais condições os incentivos à produção favorecem a abertura de novas áreas ou o aproveitamento de terras já utilizadas. Essa questão aproxima o debate sobre retornos privados e custos ambientais das discussões sobre preços, capacidade produtiva e instituições. Para o problema proposto neste artigo, a revisão dessas contribuições busca esclarecer os mecanismos sugeridos pela literatura, as evidências que os sustentam e os limites de sua aplicação ao contexto de Mato Grosso entre os anos de 2001 e 2024. A trajetória econômica e ambiental do estado oferece o contexto para discutir a pertinência dessas explicações e suas possíveis diferenças entre atividades agropecuárias e biomas.

% Origem: introdução, REF1--REF2 e fundamentação conceitual de REF7/REF11.
\subsection{Retornos privados e custos ambientais da expansão agropecuária}
\label{literatura-retornos}

A expansão agropecuária envolve uma escolha econômica entre usos alternativos da terra. A conversão da vegetação em área produtiva para lavoura ou pecuária pode ampliar a renda privada associada à produção, ao mesmo tempo em que reduz estoques de capital natural e serviços ecossistêmicos que nem sempre são incorporados aos retornos percebidos pelos agentes. Essa diferença entre o valor social da vegetação e os benefícios que o proprietário consegue apropriar integra o enquadramento de \citeonline{Angelsen2010}, que distingue as rendas da agricultura e da floresta. A decisão depende, portanto, da comparação entre alternativas efetivamente acessíveis ao agente, e não apenas do valor da produção que poderá ser obtida na área convertida.

Os efeitos da conversão podem ultrapassar os limites da propriedade convertida, embora sua extensão dependa dos processos considerados. No Cerrado brasileiro, por exemplo, \citeonline{Schwamback2024} documentaram maior escoamento de nutrientes sob pastagem e maior perda de solo sob cana, em comparação com o cerrado arborizado. Na fronteira Amazônia-Cerrado, a substituição de floresta por soja esteve associada à redução do carbono do solo, com respostas diferentes conforme a profundidade e o manejo \cite{Popin2025}. Esses estudos documentaram alterações biofísicas, mas não mensuraram, por si sós, o custo social total da conversão.

Além das alterações biofísicas documentadas, a conversão pode ter implicações econômicas ao modificar serviços do solo que sustentam a produção ao longo do tempo. Nessa perspectiva, \citeonline{Brady2015} propuseram valorar esses serviços a partir de seus efeitos sobre fluxos futuros de renda agrícola, utilizando dados de experimentos de longo prazo na Europa e nos Estados Unidos que compararam a produtividade agrícola sob diferentes práticas de adubação e manejo do solo.

A comparação entre benefícios e custos privados, como o lucro de um fazendeiro, e públicos, como os impactos gerados para outros agentes e territórios, enfrenta dificuldades adicionais. Os serviços ambientais diferem entre lugares, operam em escalas distintas e muitos não possuem preços diretamente observáveis. Estudos que integram informações econômicas e biofísicas mostram a importância da localização e das alternativas de uso da terra para valorar esses serviços e orientar a conservação \cite{Crossman2009,Gascoigne2011}. Assim, a oposição entre conservar e produzir pode refletir diferenças entre retornos apropriados privadamente e custos distribuídos socialmente.

Para a Amazônia brasileira, \citeonline{Mann2012} estimaram esse descompasso ao comparar o retorno econômico potencial da conversão da floresta para a produção de soja com o valor monetário dos serviços ecossistêmicos mantidos pela vegetação nativa. Para tanto, combinaram estimativas do retorno da soja calculadas por um modelo de produtividade e pela diferença entre receitas e custos, com a valoração dos serviços ecossistêmicos por preços-sombra derivados de uma função de produção Cobb-Douglas, que incorporou a produtividade primária da vegetação como indicador desses serviços. Sob as hipóteses adotadas, a conclusão foi que o valor dos serviços superava a renda da soja em parcela expressiva da área analisada, e a omissão desses serviços reduzia o bem-estar social relativamente ao ótimo modelado.

No caso do Pantanal, existem tentativas de mensurar economicamente os serviços ambientais, mas os próprios resultados evidenciam dificuldades de escala, cobertura e comparabilidade. \citeonline[p.~3]{SeidlMoraes2000}, ao recalcularem para a sub-região da Nhecolândia\footnote{Região do Pantanal sul-mato-grossense (Bacia do Alto Paraguai) caracterizada por florestas sazonais, pastagens nativas e lagoas alcalinas, com tradicional manejo da pecuária bovina extensiva.} estimativas globais com dados locais, obtiveram valor anual de aproximadamente US\$ 5,8 mil por hectare, em dólares de 1994, equivalente a cerca de metade daquele derivado dos parâmetros globais, diferença atribuída às condições biofísicas regionais.

Mais recentemente, \citeonline[p.~7]{FairbrassEtAl2025}, ao estruturarem indicadores de capital natural para o Pantanal, buscaram reunir informações sobre a extensão, o estado de conservação e o valor econômico dos ecossistemas, os serviços ambientais oferecidos e as atividades humanas que deles se beneficiam. Dos 29 indicadores previstos para os ativos ecossistêmicos, encontraram dados que atendiam aos critérios do estudo para apenas 14, menos da metade do conjunto.

Essa cobertura incompleta limitou a caracterização das relações entre a economia e o ambiente. Embora houvesse informações sobre atividades com valor de mercado, como pecuária, soja, pesca e turismo, faltavam indicadores sobre a condição e o valor econômico dos ecossistemas que as sustentam, sobre parte das emissões geradas por essas atividades e acerca dos processos ecológicos essenciais, como a ciclagem de nutrientes e a manutenção da biodiversidade.

As dificuldades de cobertura observadas para o Pantanal se estendem a outros biomas. \citeonline{CampoliStivali2026} reuniram estudos de valoração econômica da vegetação nativa nos seis biomas brasileiros (Amazônia, Cerrado, Pantanal, Mata Atlântica, Caatinga e Pampa) e padronizaram estimativas de custo social do desmatamento em dólares de 2025 por hectare. Os autores concluíram que os estudos diferiam tanto em relação aos benefícios da vegetação incluídos na valoração econômica quanto nos métodos empregados.

As diferenças entre os valores econômicos estimados para os serviços ambientais podem refletir tanto as características dos biomas como as diferenças no que cada estudo conseguiu mensurar. Em conjunto, os trabalhos discutidos mostram por que o retorno privado da conversão não representa, necessariamente, seu resultado econômico para a sociedade: a renda obtida com o novo uso da terra pode coexistir com perdas de serviços ambientais cujos custos não recaem integralmente sobre quem decide converter.

Essa distinção fundamenta o recorte deste artigo, voltado às decisões que transformam o estoque de vegetação nativa. Busca-se investigar como os incentivos privados e as condições institucionais se relacionam com a conversão para diferentes usos agropecuários, sem supor que essa investigação resolva as lacunas de valoração. Cabe, portanto, examinar como os preços agropecuários e a capacidade produtiva da terra modificam a atratividade das alternativas disponíveis aos produtores.

% Origem: REF3, REF6--REF8 e diálogo empírico de REF11.
\subsection{Preços agropecuários e decisões de uso da terra}
\label{literatura-precos}

No argumento deste artigo, os preços dos produtos agropecuários e a capacidade produtiva da terra representam dimensões observáveis dos incentivos ao uso agropecuário. Mantidos os custos e as demais condições, uma elevação de preço representa maior acréscimo de receita por hectare onde a produção potencial do uso beneficiado é maior. Essa relação econômica ajuda a fundamentar a hipótese apresentada na introdução. Sua tradução em conversão efetiva, contudo, depende da vegetação disponível, dos custos de abertura da área, das alternativas em terras já convertidas e das restrições institucionais. Trata-se de uma previsão condicional a investigar, não de uma consequência necessária de qualquer aumento de produtividade.

A comparação dos retornos relativos orienta a análise das políticas de redução do desmatamento. \citeonline{Angelsen2010} organiza essas políticas em torno das rendas relativas da agricultura e da floresta, distinguindo intervenções que reduzem a atratividade da expansão, elevam os benefícios apropriáveis da conservação ou restringem diretamente a conversão. Esse enquadramento permite compreender por que a intensificação produtiva pode ampliar a pressão local de expansão em determinadas condições, sem presumir que produzir mais implique necessariamente converter mais terra.

A literatura empírica consolida o papel dos atributos biofísicos e locacionais nas decisões de uso do solo. Para os Estados Unidos, \citeonline{LubowskiEtAl2008} demonstraram que os proprietários respondem aos retornos líquidos dos usos alternativos em interação com a qualidade da terra. Na Amazônia brasileira, \citeonline{Pfaff1999} evidenciou a importância das características da terra e dos custos de transporte para o desmatamento, advertindo que a alocação espacial de políticas públicas também responde a esses atributos. Em escala espacial detalhada, \citeonline{ChomitzThomas2003} confirmaram que precipitação, solos, acesso a mercados e proteção territorial condicionam a conversão florestal e a intensidade da pecuária. Contudo, documentar que terras mais férteis ou acessíveis concentram maior ocupação agropecuária não equivale, por si só, a demonstrar que esses territórios respondem com maior intensidade a choques de preços, relação interativa que constitui o foco da hipótese deste artigo.

Testar essa sensibilidade diferencial aos preços exige, ademais, distinguir conceitualmente a capacidade produtiva predeterminada da produtividade observada. O rendimento agrícola aferido em áreas já cultivadas (como a produtividade média municipal) é endógeno às decisões de manejo, ao uso de insumos e à própria seleção das parcelas já abertas. Por essa razão, avaliar como a aptidão do território amplifica os choques de preços requer medidas de capacidade produtiva anteriores à decisão de converter. A combinação entre preços e capacidade produtiva capta variações na receita bruta potencial, funcionando como uma \textit{proxy} observável dos incentivos de mercado, e não como uma medida direta do retorno líquido esperado, que depende de custos de conversão, expectativas e restrições de crédito.

Em escala tropical, \citeonline{BermanEtAl2023} combinaram preços internacionais e aptidão agrícola para examinar a relação entre incentivos de mercado e desmatamento entre 2001 e 2018.\footnote{A pesquisa abrangeu regiões tropicais de diferentes países, com destaque para a Amazônia, o Sudeste Asiático e a Bacia do Congo, onde a associação entre preços agrícolas e desmatamento foi mais forte.} Os autores encontraram uma associação positiva entre preços agrícolas e desmatamento, cuja intensidade variou conforme as características econômicas locais. Essa abordagem é pertinente porque permite considerar que um mesmo aumento de preço pode gerar incentivos distintos à conversão, dependendo das possibilidades produtivas de cada território.

O uso da aptidão agrícola para medir a exposição aos preços, contudo, não demonstra, por si só, que a capacidade produtiva local modifica a resposta da conversão da forma prevista nesta pesquisa. Há um desafio de identificação: os preços internacionais podem refletir choques de oferta, sendo necessário justificar as condições sob as quais a exposição a suas variações pode ser tratada como exógena ao desmatamento analisado.

A relação entre preços e conversão depende do uso que se expande e daquele que é substituído. Por exemplo, na bacia do Tapajós, \citeonline{Braganca2018} encontrou que a elevação do preço relativo de lavouras em relação à carne favoreceu a substituição de pastagens por cultivos (\textit{pasture-to-cropland}) e reduziu o desmatamento local.

À luz desses antecedentes, a contribuição aqui pretendida consiste em examinar as condições territoriais em que os incentivos de preços se associam à conversão da vegetação nativa, avançando sobre a relação documentada por \citeonline{BermanEtAl2023}, \citeonline{AssuncaoEtAl2015} e \citeonline{HardingEtAl2021}, dentre outros. O diferencial do artigo está em analisar conjuntamente a capacidade produtiva local, o ambiente institucional e os destinos da conversão, distinguindo a abertura de áreas para agricultura e pastagem da reorganização de usos em áreas anteriormente convertidas, destacada por \citeonline{Braganca2018}.

Para Mato Grosso, busca-se identificar quais atividades e territórios apresentam maior sensibilidade aos preços e em que medida essa resposta varia com as condições produtivas e institucionais. A análise poderá, assim, oferecer evidências sobre o alcance e os limites das explicações existentes e subsídios à discussão de políticas de conservação que considerem essas diferenças. Poderá abrir perspectivas para que estudos futuros investiguem a adequação territorial de instrumentos regulatórios e econômicos de conservação em outros contextos, regiões e biomas.

% Origem: REF1--REF2 e REF9; limites de escopo estabelecidos ao final de REF11.
\subsection{Ambiente institucional e respostas aos incentivos econômicos}
\label{literatura-instituicoes}

As instituições ajudam a compreender por que os produtores podem responder de formas distintas aos incentivos de preços discutidos anteriormente. Na economia institucional, \citeonline{North1991} compreende as instituições como regras formais, a exemplo de leis, regulamentos e direitos de propriedade, e restrições informais, como costumes e normas de conduta compartilhados por membros de uma comunidade. Igualmente, \citeonline{North1993} distingue essas regras das organizações que atuam sob elas e participam de sua elaboração e aplicação, entre as quais órgãos públicos, empresas e associações.

No que se refere à governança dos recursos naturais, \citeonline{Ostrom2009} destaca a relevância das regras efetivamente praticadas, do monitoramento e das sanções, inclusive quando organizados pelos próprios usuários. Por seu turno, \citeonline{Shimshack2014} examina como o monitoramento e a aplicação das regras influenciam o cumprimento das exigências ambientais. Essas contribuições permitem situar as respostas aos incentivos na articulação entre normas, organizações e condições de implementação.

Assim sendo, o ambiente institucional compreende as regras formais e informais que orientam as decisões de conversão, conservação e uso dos recursos ambientais, bem como os mecanismos que sustentam sua aplicação. Inclui leis e regulamentos sobre o uso da terra e da vegetação, direitos de propriedade, costumes e normas de conduta reconhecidos pelos agentes. Sua implementação envolve organizações públicas e privadas e os próprios usuários dos recursos, por meio de monitoramento, fiscalização e aplicação de sanções legais, contratuais ou sociais. Nos termos da hipótese suscitada neste trabalho, importa compreender como essa articulação modifica os custos, os riscos e os retornos esperados das alternativas de uso da terra, distinguindo a existência de uma regra das condições que tornam seu cumprimento efetivo.

Sob essa perspectiva, o enquadramento econômico não supõe efeitos uniformes das instituições. \citeonline{PachecoMeyer2022} compararam regimes fundiários em diferentes contextos brasileiros e mostraram que os resultados dependeram do regime de origem, da alternativa de destinação da terra e do contexto socioambiental. A comparação de terras privadas com terras públicas sem destinação não produz a mesma conclusão que sua comparação com áreas protegidas.

Já \citeonline{ProbstEtAl2020}, ao examinarem o programa Terra Legal \footnote{O Programa Terra Legal foi uma iniciativa do governo federal brasileiro, criada em 2009 (Lei nº 11.952/2009), voltada para a regularização fundiária de ocupações informais em terras públicas federais na Amazônia Legal. O programa priorizava a doação ou venda facilitada de terras a pequenos e médios agricultores, atrelando a concessão do título definitivo ao cumprimento de contrapartidas ambientais previstas no Código Florestal.}, encontraram aumento do desmatamento entre pequenos e médios proprietários titulados e maior resposta aos preços agropecuários. A regularização, desse modo, não deve ser confundida com um indicador inequívoco de conservação.

Esses resultados situam os direitos de propriedade entre as condições que podem modificar os retornos esperados do uso da terra. Para o argumento deste artigo, importa compreender quais alternativas de produção e conservação são favorecidas por diferentes regras de acesso e uso, considerando as diferenças entre regimes fundiários e iniciativas de titulação. A investigação específica da apropriação irregular e da valorização especulativa da terra permanece, conforme delimitado na introdução, reservada a outro artigo.

Além das regras de acesso à terra, instrumentos de registro ambiental podem modificar as condições em que os produtores tomam suas decisões. Ao analisarem o Cadastro Ambiental Rural (CAR) no Pará e em Mato Grosso, \citeonline{JungEtAl2022} empregaram um modelo dinâmico com procedimentos para mitigar a seleção no cadastro. Os autores não encontraram efeito estatisticamente significativo da inscrição sobre a conversão para lavouras, mas observaram expansão de pastagens, com diferenças associadas ao porte dos estabelecimentos e à conformidade ambiental. O resultado para lavouras não demonstra que o efeito seja exatamente zero.

O estudo também encontrou evidências de que mudanças no acesso ao crédito podem ser um dos mecanismos associados à expansão de pastagens. O cadastro precisa, assim, ser considerado tanto em sua função ambiental quanto em sua relação com as oportunidades de financiamento da produção. Essa articulação ajuda a compreender por que a existência de uma exigência ambiental não basta para antecipar seus efeitos sobre o uso da terra, que dependem dos incentivos econômicos associados à sua implementação.

No âmbito privado, a Moratória da Soja, examinada por \citeonline{HardingEtAl2021}, exemplifica a implementação de critérios de conservação por meio das relações de compra na cadeia produtiva. Esse caso situa os critérios comerciais entre os mecanismos que podem alterar as alternativas dos produtores. Instrumentos econômicos de conservação, como pagamentos por serviços ambientais, mercados de carbono ou financiamento ambiental, integram esse contexto. As receitas potenciais associadas à conservação do carbono florestal, por exemplo, podem modificar os retornos relativos enfrentados pelos agentes \cite{Angelsen2010,SticklerEtAl2013}.

Em Mato Grosso, esse enquadramento orienta o exame da sensibilidade da conversão aos preços em contextos institucionais distintos, considerando a capacidade produtiva local e os destinos agropecuários da área convertida. Para essa análise, os registros de autuações do Instituto Brasileiro do Meio Ambiente e dos Recursos Naturais Renováveis (IBAMA) constituem uma fonte candidata para observar uma dimensão da atuação fiscalizatória. Esses registros não representam, isoladamente, a efetividade das regras nem permitem mensurar todo o ambiente institucional, incluindo costumes e normas informais. A análise pretende relacionar as dimensões institucionais que os dados permitirem observar às diferenças na resposta da conversão, mantendo a definição conceitual mais ampla como contexto de interpretação. Cabe, a seguir, situar essas relações na trajetória econômica e ambiental de Mato Grosso desde os anos 2000.

% Origem: REF2, REF9 e comparação entre instrumentos em REF11.
% Origem: REF4--REF6, REF10 e contribuição/transição empírica de REF11.
\subsection{Síntese da literatura e contexto de Mato Grosso}
\label{literatura-sintese}

A situação econômica e ambiental de Mato Grosso foi moldada pela expansão agropecuária, pela integração aos mercados internacionais e pelas mudanças nas regras de uso da terra. Desde os anos 2000, o crescimento produtivo combinou a incorporação de terras, investimentos de capital e avanços tecnológicos, enquanto a proteção da vegetação passou a envolver diferentes instrumentos públicos e privados. Nessa trajetória articulam-se os incentivos econômicos e as condições institucionais discutidos anteriormente.

Isto posto, a inserção nos mercados internacionais aparece na reorientação das exportações de soja. No contexto posterior à adesão chinesa à Organização Mundial do Comércio (OMC), \citeonline{LathuilliereEtAl2014} documentaram que, entre 2006 e 2010, a China superou a Europa como destino das exportações de Mato Grosso, acompanhando mudanças na escala e nas implicações ambientais da produção. Essa reorientação compõe o contexto de demanda externa no qual se desenvolveram as atividades agropecuárias do estado.

Entre 2000 e 2024, a produção de soja em Mato Grosso passou de aproximadamente 8,8 para 38,4 milhões de toneladas, enquanto a de milho aumentou de 1,4 para 47,4 milhões de toneladas \cite{IBGEPAM2026}. No mesmo intervalo, o rebanho bovino passou de 18,9 para 32,9 milhões de cabeças \cite{IBGEPPM2026}.

Esse crescimento ocorreu em meio a mudanças tecnológicas que ampliaram as possibilidades de produção agropecuária em Mato Grosso, inclusive em terras já utilizadas. Entre essas mudanças, \citeonline{Spera2017} destacou a adaptação de cultivares de soja à agricultura tropical, a correção dos solos e o desenvolvimento de variedades de ciclo curto, que viabilizam o cultivo sequencial de soja e milho no mesmo ano agrícola. Complementando, \citeonline{VanWeyEtAl2013} mostraram que a expansão de duas safras esteve associada ao acesso a transportes, à produção agrícola anterior e à conexão com mercados.

O aproveitamento de áreas já utilizadas aparece principalmente na trajetória de expansão da soja. Na porção florestal de Mato Grosso, \citeonline{MacedoEtAl2012} verificaram que 91\% da expansão dessa cultura entre 2006 e 2010 ocorreu em terras anteriormente abertas, sobretudo pastagens, ao mesmo tempo em que a produção cresceu e o desmatamento diminuiu. A coexistência, nesse período, de crescimento produtivo e redução do desmatamento dá conteúdo à distinção entre ampliar a produção e converter vegetação, também presente na análise de substituição de usos de \citeonline{Braganca2018}.

A intensificação produtiva não assegura, por si só, a conservação, pois seus efeitos dependem das alternativas de uso da terra e dos incentivos presentes em cada contexto \cite{Angelsen2010}. Na porção florestal de Mato Grosso, o crescimento da produção em áreas já utilizadas coexistiu com novas conversões para pastagem, que responderam pela maior parte do desmatamento remanescente no período examinado por \citeonline{MacedoEtAl2012}.

Dentre as ações institucionais com impactos ambientais em Mato Grosso nesse período, o Plano de Ação para Prevenção e Controle do Desmatamento na Amazônia Legal (PPCDAm), lançado em 2004, e o uso de alertas de satélite pelo Sistema de Detecção de Desmatamento em Tempo Real (DETER) fortaleceram a articulação entre monitoramento e fiscalização \cite{AssuncaoEtAl2023}.

A Moratória da Soja, iniciada em 2006, acrescentou uma restrição comercial à expansão: empresas signatárias se comprometeram a não comprar soja produzida em áreas do bioma Amazônia desmatadas após a data de corte do acordo, posteriormente ajustada para 22 de julho de 2008. \citeonline{HeilmayrEtAl2020} estimaram redução do desmatamento associada à Moratória em áreas aptas à soja e destacaram sua complementaridade com cadastros públicos e monitoramento. A abrangência desse compromisso conferiu às áreas amazônicas do estado uma condição comercial distinta da observada no Cerrado e no Pantanal, aos quais os mesmos critérios não se estendiam.

No plano das políticas públicas, o novo Código Florestal (Lei n.~12.651/2012) reestruturou as regras de proteção da vegetação nativa e instituiu o Cadastro Ambiental Rural (CAR) como registro nacional obrigatório dos imóveis rurais. A legislação estabeleceu percentuais de Reserva Legal que variam conforme a localização e a formação vegetal, além de prever instrumentos específicos para a regularização de passivos ambientais \cite{Brasil2012}.\footnote{O art.~12 estabelece, como regra geral na Amazônia Legal, 80\% para imóveis em áreas de florestas, 35\% em áreas de cerrado e 20\% em campos gerais, observadas as exceções legais. A Amazônia Legal é um recorte territorial distinto do bioma Amazônia.}

Para Mato Grosso, \citeonline{SticklerEtAl2013} examinaram os custos e as potenciais consequências do cumprimento dessas obrigações, sublinhando os desafios impostos pela instabilidade regulatória e a importância de combiná-las com incentivos à conservação. Essa complexidade regulatória e a assimetria nos custos de adequação ajudam a compreender por que um instrumento unificado como o CAR produz respostas tão distintas entre atividades e portes de produtores, a exemplo do que documentaram \citeonline{JungEtAl2022}.

A cooperação internacional acrescentou às regras de proteção ambiental possibilidades de financiamento da conservação. O Acordo de Paris, adotado em 2015, estabeleceu um marco de cooperação climática no qual os países apresentam compromissos nacionais de redução de emissões. No campo florestal, seu art.~5 incentiva a conservação dos estoques de carbono e o apoio à redução das emissões por desmatamento e degradação, inclusive por meio de pagamentos por resultados \cite{UNFCCC2015}.

Em Mato Grosso, a articulação entre produção e conservação ganhou expressão na estratégia Produzir, Conservar e Incluir (PCI), apresentada na vigésima primeira Conferência das Partes (COP21), em 2015. Construída com a participação do governo estadual, do setor produtivo e de organizações da sociedade civil, a PCI busca conciliar o aumento da produção em áreas já abertas com a redução do desmatamento e a inclusão da agricultura familiar e de comunidades tradicionais \cite{FagundesDomingues2017}.

O financiamento internacional dessas iniciativas ganhou um instrumento concreto em 2017, quando Mato Grosso formalizou acordos com Alemanha e Reino Unido no âmbito do programa REDD Early Movers (REM). O programa remunera resultados de redução das emissões associadas ao desmatamento, direcionando recursos para ações de conservação, produção sustentável e apoio à agricultura familiar e às comunidades tradicionais. Nos acordos firmados pelo estado, os recursos seriam operacionalizados pelo Fundo Brasileiro para a Biodiversidade (Funbio) \cite{FagundesDomingues2017}. Esse arranjo vincula resultados ambientais à obtenção de recursos para políticas e projetos, criando uma forma de incentivo econômico à manutenção da vegetação e ao uso produtivo de áreas já abertas.

Além do apoio financeiro à conservação, o monitoramento, a fiscalização e as restrições comerciais podem modificar os incentivos ao uso da terra. Na Amazônia brasileira, \citeonline{AssuncaoEtAl2023}, ao explorarem limitações de detecção por nuvens no DETER, identificaram efeito de contenção do desmatamento pelo monitoramento e pela fiscalização. Em Mato Grosso, \citeonline{Sousa2016} examinou multas, embargos e apreensões, destacando diferenças entre essas sanções e dificuldades em sua aplicação. \citeonline{HardingEtAl2021} encontraram respostas distintas do desmatamento aos preços conforme as políticas e as áreas consideradas na Amazônia brasileira, incluindo resultados compatíveis com substituições de cultivos e deslocamentos da pressão para terras não protegidas.

Em Mato Grosso, a interpretação dessas relações envolve territórios com diferentes características biofísicas e trajetórias de uso da terra. \citeonline{SimoesEtAl2020} produziram mapas anuais de cobertura e uso da terra entre 2001 e 2017, classificando séries de imagens do sensor de satélite MODIS, com resolução espacial de 250 m, por aprendizado de máquina. Os mapas distinguiram floresta, vegetação de cerrado, pastagens e sistemas agrícolas de uma ou duas safras, como soja seguida de milho ou algodão. Os autores registraram expansão das lavouras e dos cultivos de duas safras, além de crescimento das pastagens sobretudo ao norte, no bioma Amazônia.

A dinâmica das pastagens ganha destaque no estudo de \citeonline{SeidlEtAl2001}, que examinou indicadores econômicos e incentivos dos proprietários relacionados à expansão de pastagens cultivadas no Pantanal brasileiro, incluindo áreas fora de Mato Grosso. Essa perspectiva amplia a discussão para além da expansão de grãos e ajuda a situar a diversidade de usos agropecuários presente no estado.

De todo modo, os desafios na governança ambiental de Mato Grosso permanecem, como evidenciam os resultados de \citeonline{AssuncaoEtAl2023} e \citeonline{HardingEtAl2021}. \textcolor{red}{CONFIRMAR - Segundo dados da Secretaria de Estado de Meio Ambiente (SEMA-MT), em 2024, 99,5\% dos imóveis rurais do estado estavam inscritos no CAR, mas apenas 1,5\% haviam regularizado suas pendências ambientais. A situação evidencia que a inscrição no CAR não garante, por si só, a conformidade ambiental, sendo necessário considerar o contexto institucional mais amplo para compreender as decisões de conversão da vegetação nativa. Além disso, do total de multas aplicadas entre 2001 e 2024, xx\% apenas foram pagas, xx\% foram parceladas e xx\% permaneciam em aberto, indicando que a aplicação de sanções também enfrenta desafios de efetividade.}

Esse quadro orienta a comparação proposta neste artigo entre a conversão bruta para agricultura e pastagem, sob uma base territorial comum. Busca-se compreender se a resposta aos preços varia com a capacidade produtiva local e com as dimensões institucionais que os dados permitem observar, situando as relações discutidas na literatura no contexto de Mato Grosso. As decisões de converter modificam o estoque de vegetação disponível nos períodos seguintes, razão pela qual ele é tratado como componente endógeno do processo econômico. O contraste entre atividades e biomas poderá, assim, esclarecer o alcance e os limites das relações propostas na hipótese do trabalho.
